\documentclass[11pt,a4paper]{article}
\usepackage[T1]{fontenc}
\usepackage{lmodern,amsmath,amssymb,amsthm,microtype}
\usepackage[margin=27mm]{geometry}
\usepackage[colorlinks=true,linkcolor=blue,urlcolor=blue]{hyperref}
\setlength{\emergencystretch}{2em}
\newtheorem{theorem}{Theorem}[section]
\newtheorem{lemma}[theorem]{Lemma}
\newtheorem{corollary}[theorem]{Corollary}
\newcommand{\Q}{\mathbb Q}\newcommand{\Z}{\mathbb Z}
\DeclareMathOperator{\Res}{Res}
\title{Finite reconstruction of elliptic Lucas points}
\author{}\date{17 November 2026}
\begin{document}\maketitle
\begin{abstract}
Squared-Lucas triples with unbounded scalar label are Zariski dense
in the threefold power of the source's restriction-of-scalars abelian
surface. The same remains true after adjoining the shared auxiliary
point and its three multiplication points, provided each triple retains
infinitely many auxiliary indices. Consequently a fixed finite
algebraic reconstruction of a discarded Pell point cannot preserve
this flexible preparation family. This is an obstruction to that
specific elimination procedure, not a lower bound on the number of
variables in a Diophantine definition.
\end{abstract}

\section{The integer-label obstruction}
Write
\[
 F(a,t,b)=t^2-a(a-2)tb+(a-2)^2b^2-1.
\]

\begin{lemma}\label{lem:affine}
For fixed integers $u,v,w,r$, with $(u,v,w)\ne(0,0,0)$,
the integer solutions of
\[
 F(s^2,t,b)=0,\qquad b\ne0,\qquad us+vt+wb=r
\]
have bounded $|s|$.
\end{lemma}
\begin{proof}
It suffices to consider $|s|\ge3$. The elementary Lucas growth lemma
gives $\log|t|>s^2-3$; see \cite{Lucas}. Suppose first $v\ne0$.
Put $a=s^2$ and $R=r-us$. Eliminating $t=(R-wb)/v$ gives
\[
 A(a)b^2+B(s)b+C(s)=0,
\]
where
\begin{align*}
 A(a)&=w^2+vw\,a(a-2)+v^2(a-2)^2,\\
 B(s)&=R[-2w-va(a-2)],\\
 C(s)&=R^2-v^2.
\end{align*}
The leading coefficient is never the zero polynomial. Indeed
\[
 A(a)=v(w+v)a^2-2v(w+2v)a+(w^2+4v^2);
\]
if its quadratic coefficient vanishes, then $w=-v$ and
$A(a)=v^2(5-2a)\ne0$. Outside finitely many values of $s$,
the integer $A(s^2)$ is nonzero and has absolute value at least one.
The elementary polynomial-root bound therefore gives
\[
 |b|\le 1+\max(|B/A|,|C/A|)=O((1+|s|)^5).
\]
The linear relation gives the same polynomial bound for $t$,
contradicting the Lucas estimate as $|s|$ tends to infinity.

If $v=0$ and $w\ne0$, then $b=(r-us)/w=O(1+|s|)$.
Applying the root bound directly to the monic quadratic $F(s^2,t,b)$
gives $|t|=O((1+|s|)^6)$, again a contradiction. If $v=w=0$,
then $u\ne0$ and the relation fixes $s$. This proves the lemma.
\end{proof}

The integrality of the cofactor $b$ is essential. For every integer
$|s|\ge3$, setting
\[
 a=s^2,\qquad b=\frac1{a-2},\qquad t=a
\]
gives $F(a,t,b)=0$, but $\log|t|=2\log|s|$.
Thus retaining only the integrality of $(a-2)b$, as provided by a
monic equation in that product, does not retain Lucas growth.

\section{The cyclic group and its closure}

Let
\[
 K=\Q(\sqrt2),\quad d=3-\sqrt2,\quad
 E:y^2=x^3-d^2x,\quad A=\Res_{K/\Q}E,
\]
and let $\gamma\in A(\Q)$ correspond to the fixed non-torsion source
point $P\in E(K)$. Set $\Gamma=\Z\gamma$.

\begin{lemma}\label{lem:densecyclic}
The cyclic subgroup $\Gamma$ is Zariski dense in $A$.
\end{lemma}
\begin{proof}
Over an algebraic closure, identify both conjugate factors with
$E_0:Y^2=X^3-X$. The resulting points $R,R'$ are independent
modulo torsion over $\operatorname{End}(E_0)=\Z[i]$;
\cite[Section 2]{Weighted} proves this directly. Briefly, the
square classes of $3-\sqrt2$ and $3+\sqrt2$ are independent in
$K(i)^\times/K(i)^{\times2}$, so an automorphism fixes $R$ and
every endomorphism while negating $R'$. Applying this automorphism
to an endomorphism relation forces each coefficient to vanish.

The identity component of the closure of a subgroup of an abelian
variety is an abelian subvariety. If the closure of $(R,R')$ were
proper in $E_0^2$, a nonzero homomorphism $E_0^2\to E_0$ would
vanish on that component. A positive multiple of $(R,R')$ belongs
to the component, giving a nonzero endomorphism relation modulo
torsion. Independence excludes this. Hence the closure is all $A$.
\end{proof}

We use Mordell--Lang in the following familiar form: the intersection
of a closed subvariety of an abelian variety with a finitely generated
subgroup is a finite union of cosets of subgroups. For a primary
reference, the stronger uniform theorem and its explicit coset
formulation are \cite[Theorem 1.1 and equation (1.1)]{GGK}; only the
ordinary finiteness statement is needed here.

\begin{lemma}\label{lem:lattices}
If $Y\subsetneq A^m$ is closed, the integer vectors
\[
 (n_1,\ldots,n_m)\in\Z^m,
 \qquad ([n_1]\gamma,\ldots,[n_m]\gamma)\in Y
\]
lie in a finite union of proper affine rational hyperplanes.
\end{lemma}
\begin{proof}
Mordell--Lang gives finitely many cosets of subgroups of
$\Gamma^m\cong\Z^m$. A subgroup of rank $m$ has finite index.
Its closure is $A^m$ by Lemma~\ref{lem:densecyclic}, since it
contains $N\Gamma^m$ for some positive $N$. A coset contained in
the proper variety $Y$ therefore has rank at most $m-1$.
Each corresponding affine lattice satisfies a nontrivial affine
equation with integer coefficients, as required.
\end{proof}

\section{Density and finite reconstruction}

\begin{theorem}\label{thm:roots}
Let $\mathcal T\subset\Z^3$ consist of solutions $(s,t,b)$ of
$F(s^2,t,b)=0$ with $b\ne0$. If $|s|$ is unbounded on
$\mathcal T$, then
\[
 \{([s]\gamma,[t]\gamma,[b]\gamma):(s,t,b)\in\mathcal T\}
\]
is Zariski dense in $A^3$.
\end{theorem}
\begin{proof}
Otherwise Lemma~\ref{lem:lattices} would place all the integer
triples in finitely many proper affine hyperplanes. Lemma~\ref{lem:affine}
bounds $|s|$ on each of them, a contradiction.
\end{proof}

The theorem allows only one chosen Lucas solution for each scalar
label. No density in the set of integer labels is assumed.

\begin{theorem}\label{thm:auxiliary}
Retain $\mathcal T$ as above. For each $(s,t,b)\in\mathcal T$
let $\mathcal H_{s,t,b}\subset\Z$ be any infinite set. The points
of $A^7$ with index vectors
\[
 (s,t,b,h,hs,ht,hb),\qquad
 (s,t,b)\in\mathcal T,\quad h\in\mathcal H_{s,t,b},
\]
are Zariski dense.
\end{theorem}
\begin{proof}
If their closure were proper, Lemma~\ref{lem:lattices} would give
a finite list of nontrivial affine equations in the seven indices.
Upon substitution, each equation has the form
\[
 L_0(s,t,b)+hL_1(s,t,b)=0,
\]
where $L_0,L_1$ are affine integer polynomials and at least one is
nonzero. For a fixed triple infinitely many $h$ occur. At least one
equation in the finite list must hold for infinitely many of them;
for this equation both $L_0$ and $L_1$ vanish at the triple.
The zero set of a nonzero constant contains no triple; otherwise
Lemma~\ref{lem:affine} bounds $|s|$ on that fixed affine equation.
Taking the maximum over the finite list bounds $|s|$ on all of
$\mathcal T$, a contradiction.
\end{proof}

\begin{corollary}\label{cor:finite}
No fixed finite algebraic correspondence over the retained point
coordinates can reconstruct the discarded point $[t]\gamma$ from
$[s]\gamma,[b]\gamma$ on a family as in
Theorem~\ref{thm:roots}. Likewise no such correspondence can
delete one or more of the seven point roles in
Theorem~\ref{thm:auxiliary} while preserving its full preparation
family and leaving the others as its only point parameters.
\end{corollary}
\begin{proof}
The closure of a finite algebraic reconstruction has dimension at
most that of its retained point base. It is therefore a proper
closed subvariety of $A^3$, respectively $A^7$. The two density
theorems exclude it.
\end{proof}

For the old completeness construction, a fixed input fixes its input
point. Its paired input auxiliary is $[n]H$, an algebraic function
of $H$ when this input index $n$ is fixed. Adjoining these fixed-input
data therefore does not change Corollary~\ref{cor:finite}.
The old preparation permits arbitrarily large scalar labels and
infinitely many compatible auxiliary indices after its contact primes
have been selected; the theorem applies to that flexible family.

The conclusion has a deliberately limited scope. It does not exclude
a new construction that chooses only finitely many scalar labels for
each fixed input, imposes a new arithmetic relation restricting the
auxiliary index, replaces the growth clause, or reconstructs a
\emph{numerical} Pell root without reconstructing its elliptic point.
It does not assert that a rational-square cover has no individual
rational witnesses. In particular it provides no lower bound on the
number of variables defining the integers over the rationals.

\begin{thebibliography}{9}\raggedright
\bibitem{Lucas}\emph{An explicit polynomial in 29 variables over the
rationals}, repository manuscript, elementary Lucas growth lemma,
\nolinkurl{papers/few_variables/few_variables_20261108_v8.tex},
lines 349--397.
\bibitem{Weighted}\emph{Heights of positive weighted elliptic traces}, repository
manuscript, Section 2,
\nolinkurl{papers/trace_heights/weighted_trace_20261114_v27.tex}.
\bibitem{GGK}Z.~Gao, T.~Ge and L.~K\"uhne,
\emph{The Uniform Mordell--Lang Conjecture},
Publications Math\'ematiques de l'IH\'ES \textbf{143} (2026), 189--235,
\href{https://doi.org/10.5802/pmihes.26}{doi:10.5802/pmihes.26}.
Author manuscript:
\url{https://ziyangjeremygao.github.io/articles/UML.pdf},
Theorem 1.1 and equation (1.1).
\end{thebibliography}
\end{document}
